Kurs:Algebraische Kurven (Osnabrück 2025-2026)/Arbeitsblatt 23/latex
Kurs:Algebraische Kurven (Osnabrück 2025-2026)/Arbeitsblatt 23/latex

\setcounter{section}{23}


  \zwischenueberschrift{Übungsaufgaben}


\inputaufgabe

{}

{


Es sei $K$ ein
\definitionsverweis {Körper}{}{,}


\mathl{K[X_1 , \ldots , X_n]}{} der
\definitionsverweis {Polynomring}{}{}
über $K$ und

\mathl{{ \frac{   \partial    }{  \partial   X_1   } }}{} die
\zusatzklammer {formale} {} {}
\definitionsverweis {partielle Ableitung}{}{}
bezüglich $X_1$, also die Abbildung
\maabbeledisp {} { K[X_1 , \ldots , X_n] } { K[X_1 , \ldots , X_n]
} { f } { { \frac{   \partial   f   }{  \partial   X_1   } }
} {.}
Zeige, dass dies eine
$K$-\definitionsverweis {Derivation}{}{}
ist.


}

{} {}


\inputaufgabe

{}

{


Wir betrachten das
\definitionsverweis {maximale Ideal}{}{}


\mavergleichskettedisp

{\vergleichskette

{ {\mathfrak m}
}

{ =} { (X_1 , \ldots , X_n)
}

{ \subseteq} { K[X_1 , \ldots , X_n]
}

{ } {
}

{ } {
}

}
{}{}{}
im
\definitionsverweis {Polynomring}{}{}
über einem
\definitionsverweis {Körper}{}{}
$K$ und seine
\definitionsverweis {Potenzen}{}{}
${\mathfrak m}^d$. Zeige, dass die Monome


\mathbeddisp {X_1^{\nu_1 } \cdots X_n^{\nu_n }} {}

 {\sum_{i = 1}^n \nu_i < d} {}

 {} {} {} {,}
eine
$K$-\definitionsverweis {Basis}{}{}
des
\definitionsverweis {Restklassenringes}{}{}


\mathl{K[X_1 , \ldots , X_n] / {\mathfrak m}^d}{} bilden.


}

{} {}


\inputaufgabe

{}

{


Betrachte das Achsenkreuz


\mavergleichskette

{\vergleichskette

{ V(xy)
}

{ \subseteq }{ {\mathbb A}^{2}_{K}
}

{  }{
}

{    }{
}

{   }{
}

}
{}{}{}
und den zum Nullpunkt gehörigen
\definitionsverweis {lokalen Ring}{}{}
$R$ mit
\definitionsverweis {maximalem Ideal}{}{}
${\mathfrak m}$. Beschreibe explizit eine
$K$-\definitionsverweis {Basis}{}{}
für die
\definitionsverweis {Restklassenringe}{}{}
$R/{\mathfrak m}^n$ und bestimme die
\definitionsverweis {Dimensionen}{}{}
davon.


}

{} {}


\inputaufgabegibtloesung

{}

{


Es sei


\mavergleichskette

{\vergleichskette

{F
}

{ = }{ F_m + \cdots + F_d
}

{  }{
}

{    }{
}

{   }{
}

}
{}{}{}
die
\definitionsverweis {homogene Zerlegung}{}{}
eines Polynoms


\mavergleichskette

{\vergleichskette

{F
}

{ \in }{ K[X,Y]
}

{  }{
}

{    }{
}

{   }{
}

}
{}{}{}
mit


\mavergleichskette

{\vergleichskette

{m
}

{ \leq }{d
}

{  }{
}

{    }{
}

{   }{
}

}
{}{}{}
und es sei


\mavergleichskette

{\vergleichskette

{ {\mathfrak m}
}

{ = }{ (X,Y)
}

{  }{
}

{    }{
}

{   }{
}

}
{}{}{.}
Zeige, dass für jedes


\mavergleichskette

{\vergleichskette

{n
}

{ \geq }{m
}

{  }{
}

{    }{
}

{   }{
}

}
{}{}{}
die Multiplikationsabbildung
\maabbeledisp {} {K[X,Y] } { K[X,Y]
} { G } { FG
} {,}
einen injektiven, wohldefinierten
$K[X,Y]$-\definitionsverweis {Modulhomomorphismus}{}{}
\maabbdisp {} {K[X,Y]/ {\mathfrak m}^{n-m} } { K[X,Y]/ {\mathfrak m}^{n}
} {}
festlegt.


}

{} {}


\inputaufgabegibtloesung

{}

{


Es sei


\mavergleichskette

{\vergleichskette

{ e
}

{ \in }{ \N_+
}

{  }{
}

{    }{
}

{   }{
}

}
{}{}{}
und sei


\mavergleichskette

{\vergleichskette

{ M
}

{ \defeq }{ \{0\} \cup \N_{ \geq e}
}

{ \subseteq }{ \N
}

{    }{
}

{   }{
}

}
{}{}{.}
\aufzaehlungdrei{Bestimme

\mathl{nM_+}{} für


\mavergleichskette

{\vergleichskette

{ n
}

{ \in }{ \N_+
}

{  }{
}

{    }{
}

{   }{
}

}
{}{}{.}
}{Bestimme

\mathl{{ \# \left(   M \setminus nM_+  \right) }}{.}
}{Es sei $K$ ein
\definitionsverweis {Körper}{}{}
und setze


\mavergleichskette

{\vergleichskette

{ R
}

{ = }{ K[M]_{ {\mathfrak m} }
}

{  }{
}

{    }{
}

{   }{
}

}
{}{}{}
mit


\mavergleichskette

{\vergleichskette

{ {\mathfrak m}
}

{ = }{ (M_+)
}

{ \subseteq }{ K[M]
}

{    }{
}

{   }{
}

}
{}{}{.}
Bestimme

\mathl{\dim_{ K } \left(   R/ {\mathfrak m}^n    \right)}{.}
}


}

{} {}


\inputaufgabe

{}

{


Berechne für das durch die Erzeuger $4$ und $9$ gegebene
\definitionsverweis {numerische Monoid}{}{}


\mavergleichskette

{\vergleichskette

{ M
}

{ \subseteq }{ \N
}

{  }{
}

{    }{
}

{   }{
}

}
{}{}{}
die in den Abschätzungen von
Lemma 23.8
auftretenden Ausdrücke bis


\mavergleichskette

{\vergleichskette

{ n
}

{ \leq }{ 6
}

{  }{
}

{    }{
}

{   }{
}

}
{}{}{.}


}

{} {}


In einigen Aufgaben wird die Krull-Dimension eines kommutativen Ringes verwendet. Da wir uns hauptsächlich für Kurven interessieren, denen eindimensionale Ringe entsprechen, werden wir keine systematische Dimensionstheorie entwickeln.


Es sei $R$ ein
\definitionsverweis {kommutativer Ring}{}{.}
Eine Kette aus
\definitionsverweis {Primidealen}{}{}


\mavergleichskettedisp

{\vergleichskette

{ {\mathfrak p}_0
}

{ \subset} { {\mathfrak p}_1
}

{ \subset} { \ldots
}

{ \subset} { {\mathfrak p}_n
}

{ } {
}

}
{}{}{}
nennt man \definitionswort {Primidealkette der Länge}{} $n$
\zusatzklammer {es wird also die Anzahl der Inklusionen gezählt, nicht die Anzahl der beteiligten Primideale} {} {.}
Die \definitionswort {Dimension}{}
\zusatzklammer {oder \definitionswort {Krulldimension}{}} {} {}
von $R$ ist das
\definitionsverweis {Supremum}{}{}
über alle Längen von Primidealketten. Sie wird mit

\mathl{\operatorname{dim}  \left(    R   \right)}{} bezeichnet.


\inputaufgabe

{}

{


Es sei $R$ ein
\definitionsverweis {Hauptidealbereich}{}{,}
der kein
\definitionsverweis {Körper}{}{} sei. Zeige, dass die
\definitionsverweis {Krulldimension}{}{} von $R$ gleich eins ist.


}

{} {}


  \zwischenueberschrift{Aufgaben zum Abgeben}


\inputaufgabe

{4}

{


Berechne für das durch die Erzeuger

\mathl{5,8,11}{} gegebene Monoid


\mavergleichskette

{\vergleichskette

{ M
}

{ \subseteq }{ \N
}

{  }{
}

{    }{
}

{   }{
}

}
{}{}{}
die in den Abschätzungen von
Lemma 23.8
auftretenden Ausdrücke bis


\mavergleichskette

{\vergleichskette

{ n
}

{ \leq }{ 5
}

{  }{
}

{    }{
}

{   }{
}

}
{}{}{.}


}

{} {}


\inputaufgabe

{3}

{


Es sei


\mavergleichskette

{\vergleichskette

{ {\mathfrak a}
}

{ \subseteq }{ R
}

{  }{
}

{    }{
}

{   }{
}

}
{}{}{}
ein
\definitionsverweis {Ideal}{}{}
in einem
\definitionsverweis {kommutativen Ring}{}{,}
das in genau einem
\definitionsverweis {maximalen Ideal}{}{}
${\mathfrak m}$ als einzigem Primoberideal enthalten sei. Zeige, dass dann


\mavergleichskette

{\vergleichskette

{ R/{\mathfrak a}
}

{ \cong }{ R_{\mathfrak m}/{\mathfrak a} R_{\mathfrak m}
}

{  }{
}

{    }{
}

{   }{
}

}
{}{}{}
ist. Folgere daraus, dass für ein maximales Ideal ${\mathfrak m}$ in einem noetherschen kommutativen Ring die Isomorphie


\mavergleichskette

{\vergleichskette

{ R/{\mathfrak m}^n
}

{ \cong }{ R_{\mathfrak m}/{\mathfrak m}^n  R_{\mathfrak m}
}

{  }{
}

{    }{
}

{   }{
}

}
{}{}{}
für jedes $n$ gilt.


}

{} {}


\inputaufgabe

{5}

{


Es sei $K$ ein
\definitionsverweis {algebraisch abgeschlossener Körper}{}{}
und sei


\mavergleichskette

{\vergleichskette

{ R
}

{ = }{ K[X,Y]
}

{  }{
}

{    }{
}

{   }{
}

}
{}{}{}
der
\definitionsverweis {Polynomring}{}{}
in zwei Variablen. Zeige, dass $R$ die
\definitionsverweis {Krulldimension}{}{}
zwei besitzt.


}

{} {}


\inputaufgabe

{5}

{


Es sei $R$ ein
\definitionsverweis {noetherscher}{}{}
\definitionsverweis {kommutativer Ring}{}{.}
Zeige, dass folgende Aussagen äquivalent sind:
\aufzaehlungfuenf{$R$ hat
\definitionsverweis {Krulldimension}{}{}
$0$.
}{$R$ ist ein
\definitionsverweis {artinscher Ring}{}{.}
}{$R$ besitzt endlich viele
\definitionsverweis {Primideale}{}{,}
die alle
\definitionsverweis {maximal}{}{}
sind.
}{Es gibt eine natürliche Zahl $n$ mit


\mavergleichskette

{\vergleichskette

{ {\mathfrak m}^n
}

{ = }{ 0
}

{  }{
}

{    }{
}

{   }{
}

}
{}{}{}
für jedes maximale Ideal ${\mathfrak m}$.
}{Die
\definitionsverweis {Reduktion}{}{}
von $R$ ist ein Produkt von Körpern.
}


}

{} {}


\inputaufgabe

{3}

{


Es sei $R$ ein
\definitionsverweis {kommutativer Ring}{}{}
von endlicher
\definitionsverweis {Krulldimension}{}{}
$d$. Zeige, dass die Krulldimension des
\definitionsverweis {Polynomrings}{}{}
$R[X]$ mindestens

\mathl{d+1}{} ist.


}

{(Bemerkung: über einem noetherschen Grundring erhöht sich die Dimension beim Übergang zum Polynomring genau um eins, dies ist aber schwieriger zu beweisen.)} {}
